\begin{pastexam}[label={px:ncku-112-cce-ds-08}]{成大電通所 112 學年度 第 8 題}
（問答）

(15\%) The following puzzle is called a word ladder: In a word ladder puzzle, you must make the
change occur gradually by changing one letter at a time. At each step you must transform one word
into another word, you are not allowed to transform a word into a non-word. The following sequence
of words shows one possible solution to the problem ``Transform the word ``FOOL'' into the word
``SAGE''.''

\hspace*{1em}Example: FOOL -> POOL -> POLL -> POLE-> PALE -> SALE -> SAGE

Here, we are interested in figuring out the smallest number of transformations needed to turn the
starting word into the ending word.

(a) Construct an example graph by using these words: FOOL, POOL, POLL, POLE, PALE, SALE, SAGE, COOL,
FOUL, FOIL, FAIL, FALL, PALL, POLE, POPE, PAGE. Represent the relationships between the words as a
word ladder graph. (Hint: Think about the path where we get from the word ``FOOL'' to the word
``SAGE'').

(b) Discuss which search algorithm you can use to find an efficient path from the starting word to
the ending word, why it is more efficient compared with other. Discuss the complexity.

(c) There is one way to improve efficiency instead of using a graph to show the relationship.
Suppose that we have a huge number of \textbf{buckets}, each of them with a four-letter word on the
outside, except that one of the letters in the label has been replaced by an underscore. The labels
on the buckets are the keys in our dictionary. The value stored for that key is a list of words.
\textbf{Show the data structure and list the buckets using the following words} (don't list the
buckets containing only one item): FOOL, POOL, POLL, POLE, PALE, SALE, SAGE, COOL, FOUL, FOIL, FAIL,
FALL, PALL, POLE, POPE, PAGE
\end{pastexam}

\begin{pastsoln}
本題和成大智慧資安學程同年的第~8~題相同（三個系所共用同一份考卷），完整的詳解見歷屆試題~\ref{px:ncku-112-isec-ds-08}，這裡只列答案。

答案：(a) 15 個字（POLE 列了兩次）、19 條邊，FOOL 到 SAGE 最少 6 步，最短的路徑有 4 條，題目的範例是其中一條。
(b) 用 BFS（無權圖的最短路徑，資結書第~21.3~節），$O(|V| + |E|)$。 (c) 以「一個字母換成底線」為鍵的雜湊表，值是字的串列；
有兩個字以上的桶子共 15 個：\texttt{\_OOL}、\texttt{FO\_L}、\texttt{PO\_L}、\texttt{P\_LL}、\texttt{POL\_}、\texttt{P\_LE}、
\texttt{PO\_E}、\texttt{\_ALE}、\texttt{PA\_E}、\texttt{PAL\_}、\texttt{SA\_E}、\texttt{\_AGE}、\texttt{F\_IL}、\texttt{FA\_L}、
\texttt{\_ALL}。
\end{pastsoln}
